\section{Introduction}
\label{sec:intro}

The 2019 \emph{Rucho v.\ Common Cause} decision removed federal courts from
the business of adjudicating partisan gerrymandering, returning the question
entirely to state legislatures, state courts, and, occasionally, independent
commissions \citep{rucho2019}.  The immediate practical effect was to widen
the scope for strategic map drawing in states without independent redistricting
commissions; the longer-term intellectual effect has been to sharpen the
question of whether partisan distortion is a property of how maps are drawn or
a property of the single-member district system itself.

A growing empirical literature supports the latter view.  \citet{rodden2019why}
documents that the geographic sorting of Democratic voters into dense urban
cores produces systematic Republican over-representation even when maps are
drawn with no partisan intent whatsoever --- an observation subsequently
confirmed by \citet{chen2013unintentional} and extended to the block-group
resolution by the present research program.  If sorting is the cause and
single-member districts are the mechanism, then algorithmic redistricting can
at best neutralize the drawing process without eliminating the structural
outcome.  This observation motivates the research agenda documented in Track E:
a systematic examination of what representation would look like under six
alternative architectures.

\subsection{Why Explore Alternatives?}

Our motivation is analytical rather than advocacy.  Each alternative system
illuminates something different about the baseline.  Multi-member districts
(E.1) reveal the geographic-sorting gap by partially escaping it.  County
representation (E.2) reveals the cost of county-preservation requirements
imposed by many state constitutions.  National redistricting (E.3) quantifies
the compactness cost of federalism.  Partisan-similarity clustering (E.4)
demonstrates that algorithmic gerrymandering is technically feasible and that
the reform consensus against it rests on normative rather than technical
grounds.  Partisan-fairness seat allocation (E.5) shows what full
proportionality would require geometrically.  International applications (E.6)
test whether findings from US census geography generalize to parliamentary
systems with different boundary structures.

Together these six experiments constitute a controlled exploration of the
algorithmic redistricting design space.  By holding the core algorithm
(edge-weighted recursive bisection via METIS) constant and varying the
objective function and structural constraints, we can isolate the causal
contribution of each design dimension to observed outcomes.

\subsection{The Post-Rucho Policy Environment}

The proportionality debate has intensified on multiple fronts since 2019.
\emph{Moore v.\ Harper} (2023) rejected the independent-state-legislature
theory and confirmed that state courts may review redistricting under state
constitutional law \citep{moore2023}, creating new litigation venues but not
new substantive standards.  The \emph{Voting Rights Act} challenge in
\emph{Allen v.\ Milligan} (2023) preserved the Section 2 compactness and
majority-minority district framework \citep{allen2023}.  Louisiana
\emph{v.\ Callais} (2025) further sharpened the VRA compactness requirement
at the sub-district level.  None of these decisions addresses the structural
underrepresentation arising from geographic sorting; they regulate the process
of drawing districts within the single-member paradigm.

Reform proposals outside litigation have ranged from the \emph{Fair
Representation Act} (multi-member districts with ranked-choice voting) to
various state-level proportional allocation schemes.  None has been enacted at
the federal level.  The principal obstacle is not political will alone but
the genuine trade-offs that any alternative system must navigate: a proposal
that maximizes proportionality at the cost of compactness faces one class of
objections; a proposal that preserves counties at the cost of population
equality faces another.  Mapping these trade-offs rigorously is the
contribution of this paper.

\subsection{The Six Systems and This Paper's Scope}

We evaluate six systems in detail in companion papers E.1--E.6.  This synthesis
paper (E.0) provides:
\begin{enumerate}
  \item Concise descriptions of each system and its core design choice
        (Section~\ref{sec:systems}).
  \item A comparative Pareto analysis across five dimensions using a
        standardized 1--3 scoring rubric (Section~\ref{sec:tradeoffs}).
  \item An analysis of how the Rodden geographic-sorting constraint operates
        differently across system types (Section~\ref{sec:constraints}).
  \item A constitutional taxonomy distinguishing which alternatives are
        immediately available from which require amendment
        (Section~\ref{sec:constitutional}).
  \item A conclusion arguing that algorithmic single-member redistricting
        (the DIA) is the Pareto-optimal choice within current constitutional
        constraints (Section~\ref{sec:conclusion}).
\end{enumerate}

Our scoring framework is deliberately conservative: we report Polsby--Popper
compactness means from 50-state sweep outputs, proportionality deviation from
statewide vote shares, majority-minority district counts against VRA benchmarks,
and county-preservation rates from census TIGER data.  All figures are derived
from the pipeline documented in \citet{dellaLibera2026mec}.
